%% 02805 Social Graphs and Interactions -- final project paper (DRAFT generated with the revised analysis).
%% Upload main.tex, references.bib and the three fig*.pdf files to the Overleaf PNAS template.
%% Every number below comes from pipeline_revised/CONCLUSIONS.md. If the analysis is re-run, update them.
%% Red [TODO] marks are things only the group can fill in. Remove "lineno" before hand-in.
\documentclass[9pt,twocolumn,twoside,lineno]{pnas-new}
\templatetype{pnasresearcharticle}
\usepackage{xcolor}
\usepackage{booktabs}
% Float placement: let several figures share a page (2-column layout)
\setcounter{dbltopnumber}{2}\setcounter{topnumber}{3}\setcounter{totalnumber}{6}
\renewcommand{\dbltopfraction}{0.85}\renewcommand{\topfraction}{0.85}\renewcommand{\textfraction}{0.1}\renewcommand{\floatpagefraction}{0.8}\renewcommand{\dblfloatpagefraction}{0.8}
\newcommand{\todo}[1]{\textcolor{red}{[TODO: #1]}}

\title{Less linked at equal impact: women scientists in Wikipedia's network of biographies}

\author[a]{Carlota \todo{Surname}}
\author[a]{\todo{Group member 2}}
\author[a]{\todo{Group member 3}}
\affil[a]{Department of Applied Mathematics and Computer Science (DTU Compute), Technical University of Denmark, 2800 Kgs.\ Lyngby, Denmark}
\leadauthor{\todo{Lead author surname}}

\significancestatement{Wikipedia is where most people first read about a scientist, and its pages point to each other. We compared women and men scientists with the same number of citations, the same field and the same decade of birth. Women's pages receive about 40\% fewer links from other scientists' pages, link out less themselves, and are more often connected to no one. Their pages also read differently even after names and gender words are hidden. Being cited equally in science does not mean being connected equally in the encyclopedia.}

\authorcontributions{\todo{Specific contributions: who collected which data, who built the network analysis and nulls, who built the text analysis, who hand-labelled the validation samples, who made each figure, who wrote each section. ``All contributed equally'' is not accepted.}}
\authordeclaration{The authors declare no competing interests.}
\correspondingauthor{\textsuperscript{1}To whom correspondence should be addressed. E-mail: \todo{sXXXXXX@dtu.dk}}
\keywords{Wikipedia $|$ gender gap $|$ network science $|$ text analysis}

\begin{abstract}
Are women scientists as connected on Wikipedia as men with the same scientific impact? We built the hyperlink network of 9,680 English Wikipedia biographies of physicists and chemists born 1920--1980 (10,563 links) and matched 2,926 of them (606 women) to OpenAlex citation records through identifiers. At equal citations, field and birth decade, women's pages receive about 40\% fewer links from other scientists' pages (negative binomial rate ratio 0.60, 95\% CI 0.51--0.70). They also give fewer links and are isolated more often, so they are less embedded in both directions. The gap exceeds 1,000 within-stratum gender shuffles ($z=-4.6$), holds when only links in running prose are counted, and is the same size in biology. With names and gender words masked, a cross-validated classifier still separates women's from men's pages within field and decade (AUC 0.85; shuffled 0.50). The comparison covers only the 30\% of pages with identifiers, and page age and authorship are not controlled.
\end{abstract}

\dates{This manuscript was compiled on \today}
\doi{02805 Social Graphs and Interactions, DTU, Fall 2026}

\begin{document}
\maketitle
\thispagestyle{firststyle}
\ifthenelse{\boolean{shortarticle}}{\ifthenelse{\boolean{singlecolumn}}{\abscontentformatted}{\abscontent}}{}

\section*{Introduction}

\dropcap{W}ikipedia biographies are connected by hyperlinks. A link from one scientist's page to another's is a small editorial statement that the second person belongs in the first one's story: an advisor, a collaborator, a rival, a co-laureate. Pages that receive such links are easier to reach and are woven into the written history of a field. We ask whether women scientists are woven in as much as men who have had the same scientific impact.

Gender differences on Wikipedia are well documented. Women are a minority of biographies, their pages differ from men's in link structure and in wording, with more words about family and relationships, and links between biographies are asymmetric by gender \cite{wagner2015man,graells2015first,wagner2016glass}. Women who do have a page tend to be more notable than men who do, which suggests a higher bar for inclusion \cite{wagner2016glass,adams2019notable}, and women's biographies are challenged as not notable more often \cite{tripodi2023ms}. Field et al.\ showed that many text differences shrink once biographies are compared with matched ones, so what is compared with what matters \cite{field2022controlled}.

What these studies do not have is a measure of what the person did that is independent of Wikipedia. For scientists there is one: citations. Citations carry their own gender gap \cite{lariviere2013bibliometrics}, so comparing at equal citations asks a narrow question: on top of whatever inequality science already produced, does the encyclopedia add more?

Here we combine the link network of scientists' biographies with OpenAlex citation records \cite{priem2022openalex} and with the text of the pages. We find that, at equal citations, field and birth decade, women's pages receive about 40\% fewer links from other scientists' pages, are less connected in both directions, and are written in a recognisably different way.

\section*{Results}

% ---------------------------------------------------------------- RQ1
\subsection*{Women's pages receive fewer links at every citation level}
The link network is sparse: about one link per page, and 46\% of pages have no link to or from another scientist. In the comparison set, women's pages receive 0.56 links on average and men's 1.51, and 51\% of women's pages are isolates against 34\% of men's. These raw numbers are not at equal impact, because the women in the set are younger and less cited. The gap is nevertheless present inside every citation quintile (Fig.~\ref{fig:rq1quint}).

\begin{figure}%[tbhp]
\centering
\includegraphics[width=.8\linewidth]{fig1a_quintiles}
\caption{\textbf{Women's pages receive fewer links at every citation level.} Mean number of links received from other scientists' pages by citation quintile; bars are 95\% bootstrap intervals (606 women, 2,320 men).}
\label{fig:rq1quint}
\end{figure}

To compare like with like, we shuffled gender labels within cells of field $\times$ birth decade $\times$ citation quintile on the fixed network (Methods). Women's pages receive 0.66 times the links expected under shuffling ($z=-4.6$), give 0.75 times as many ($z=-4.8$), and are isolates 1.20 times as often ($z=+5.3$; Fig.~\ref{fig:rq1shuffle}). Count models agree: a negative binomial regression gives a women/men rate ratio of 0.60 (95\% CI 0.51--0.70), and a hurdle model shows that the gap has two parts. Women's pages are 9.2 percentage points less likely to receive any link (CI 4.6--13.8), and those that receive any receive 0.52 times as many (CI 0.40--0.69).

\begin{figure}%[tbhp]
\centering
\includegraphics[width=.8\linewidth]{fig1b_shuffle}
\caption{\textbf{Women's pages are less embedded in both directions.} Women's observed value divided by the value expected when gender is shuffled within field $\times$ birth decade $\times$ citation quintile (orange). Grey bars cover 95\% of 1,000 shuffles. Values left of 1 mean fewer links; values right of 1 mean more pages without links.}
\label{fig:rq1shuffle}
\end{figure}

The gap does not depend on how strata or links are defined. With OpenAlex subfield in the strata, the ratio for links received is 0.72 ($z=-3.4$). Further checks are in Fig.~\ref{fig:rq1models}. Counting only links in running prose makes the gap larger (0.48, CI 0.38--0.59), so links generated by infoboxes and lists do not create it. Estimating field $\times$ gender in one model over physics, chemistry and biology gives ratios of 0.59, 0.69 and 0.64, with no evidence that they differ ($p=0.66$). Page length does not account for the gap (0.62 with it in the model). Adding the page's own outgoing links moves the ratio to 0.81 (CI 0.72--0.92): women's pages are less embedded in both directions.

\begin{figure}%[tbhp]
\centering
\includegraphics[width=.8\linewidth]{fig1c_models}
\caption{\textbf{About 40\% fewer links at equal citations, in every field.} Women/men rate ratios for links received, with 95\% intervals: count models and prose-only links, negative binomial models with controls added one at a time, and field $\times$ gender estimated in one model on the three-field network. The dotted line marks no gap.}
\label{fig:rq1models}
\end{figure}

% ---------------------------------------------------------------- RQ2
\subsection*{Masked text still tells women's pages from men's}
After masking names, pronouns, gender words and kinship terms (Methods), a cross-validated logistic regression on TF-IDF features separates women's from men's pages with AUC 0.85 when only women and men of the same field and birth decade are compared (Fig.~\ref{fig:rq2auc}). Repeating the whole procedure 200 times with gender shuffled within field and decade gives 0.50 (95\% range 0.46--0.53), so the difference is not produced by field or cohort composition.

\begin{figure}%[tbhp]
\centering
\includegraphics[width=.8\linewidth]{fig2a_auc}
\caption{\textbf{With names and gender words masked, the text still separates women's from men's pages.} Cross-validated AUC for predicting gender from masked text, counting only woman--man pairs of the same field and birth decade (orange line), against 200 repeats of the full procedure with gender shuffled within those cells (grey). The dotted line marks chance.}
\label{fig:rq2auc}
\end{figure}

The words behind this, ranked by log-odds within field and birth decade, are about training and career steps on women's pages (\emph{education, career, earned, completed, doctorate}) and about the work and its standing on men's (\emph{theory, quantum, papers, known}; Fig.~\ref{fig:rq2words}). The ranking is stable: words selected on one random half of the pages point the same way in the other half 93\% of the time. Gender is also stated outright far more often on women's pages: in the unmasked text, 15.5\% of women's pages and 0.1\% of men's state the person's gender (``first woman to'').

\begin{figure}%[tbhp]
\centering
\includegraphics[width=.8\linewidth]{fig2b_words}
\caption{\textbf{Women's pages describe training and career steps; men's pages describe the work and its standing.} The 16 words most typical of women's pages (left) and of men's pages (right): absolute $z$-score of the log-odds that a page uses the word, combined within field $\times$ birth decade. Black tick marks show the value without the birth-decade adjustment.}
\label{fig:rq2words}
\end{figure}

% ---------------------------------------------------------------- RQ3
\subsection*{Links to women's pages come disproportionately from women's pages}
Of the 696 links that point to women's pages, 154 (22.1\%) come from women's pages, although women's pages are 9.6\% of all linking pages. Shuffling gender within field and birth decade predicts 16.1\%, so field and cohort explain part of this excess but not all of it.

Women could also be concentrated in particular topic clusters. We therefore shuffled gender within Louvain link community \cite{blondel2008fast} $\times$ birth decade, over ten partitions (Methods). This null predicts 13.1--13.7\%, well below the observed share ($z$ between 6.2 and 7.3; Fig.~\ref{fig:rq3null}). Concentration does not explain it either: women are present in all 16 communities with at least 100 pages, at shares from 5\% to 18\% (Fig.~\ref{fig:rq3comm}). The few links women's pages receive therefore depend more on other women's pages than topic and cohort explain.

\begin{figure}%[tbhp]
\centering
\includegraphics[width=.8\linewidth]{fig3a_homophily}
\caption{\textbf{Same-gender linking survives a null that respects topic clusters.} Share of links to women's pages that come from women's pages (orange line) against 1,000 gender shuffles within field $\times$ birth decade (light grey) and within link community $\times$ birth decade (dark grey).}
\label{fig:rq3null}
\end{figure}

\begin{figure}%[tbhp]
\centering
\includegraphics[width=.8\linewidth]{fig3b_communities}
\caption{\textbf{Women are present in every large link community, unevenly.} Share of women in each Louvain link community with at least 100 pages, labelled by its most central page; dot size is community size. The dotted line is the share over all these communities.}
\label{fig:rq3comm}
\end{figure}

\section*{Discussion}

At equal citations, field and birth decade, women scientists' Wikipedia pages receive about 40\% fewer links from other scientists' pages, link out less, and are more often connected to no other scientist.

This does not mean readers find these pages less: we measured links between biographies, not page views or links from topic pages. It is not a causal estimate of what editors do because of a person's gender. And it describes the 30\% of pages with an ORCID or OpenAlex identifier, among whom women are younger than men (33\% vs.\ 11\% born 1970 or later). Without the citation control, the gap is the same size among the 6,747 pages we could not match (rate ratio 0.50) as among matched ones (0.51), so the matching does not create it. To test it at equal citations we matched a random sample of 601 unmatched people to OpenAlex by name, accepting only unambiguous matches (82 women, 106 men): the rate ratio is 0.63 (CI 0.38--1.05), close to the 0.60 of the same model on the identifier-matched set but too imprecise to confirm it on its own.

A limitation we found ourselves changed the headline. Our first analysis used the logarithm of PageRank and found a 14\% gap. Because almost half of the pages sit at the PageRank floor, that scale compresses the difference, and it also made biology look like an exception. On the count scale the gap is about 40\% and biology is like the other fields. A second one concerns the text. Among the words that identify women's pages are British spellings and a recognisable biography style. That is what we would see if many women's pages were written recently by few editors, in which case the classifier partly detects who wrote the page. Page creation dates and creators would test this and the page-age explanation for the link gap; we did not collect them. Finally, both rounds of hand labels (100 OpenAlex matches, 300 sentences) were done with AI assistance. A second round agreed closely with the first (all 100 matches correct in both; Cohen's $\kappa=0.99$ for personal and $0.98$ for professional sentences), which shows the labels are stable but is not a check by a person.

Earlier work found gender asymmetries in how biographies link and read \cite{wagner2015man,wagner2016glass,graells2015first}. Our contribution is to show that for scientists the link asymmetry remains after an external measure of impact is held equal, and that it is a matter of embeddedness in both directions. The next study would add page creation date and creator, and name-match a sample of people without identifiers.

\section*{Materials and Methods}

\subsection*{Data collection}
On 13 November 2026 we queried Wikidata \cite{vrandecic2014wikidata} for humans with occupation physicist, chemist or biologist, born 1920--1980, with an English Wikipedia page (12,773 people). Wikitext was downloaded on 11 and 13 November 2026 with pauses between requests. Citations, works and h-index came from OpenAlex \cite{priem2022openalex}, matched only through the OpenAlex ID or ORCID stored on Wikidata; citations were summed over split author profiles. Every query, date and record count is in the data log. Pages with fewer than 50 words of prose, redirects and disambiguation pages were dropped (12,463 pages kept; 9,680 in physics and chemistry). Matches with name similarity below 0.6 were dropped. The comparison set is the 2,926 physics and chemistry pages with a match, a binary gender on Wikidata and at least one citation. For the selection check, a random sample of 300 unmatched women and 301 unmatched men, stratified by birth decade, was searched by name in OpenAlex on 18 November 2026; a profile was accepted only if it was the single candidate with at least 10 works, name similarity of at least 0.92 and a natural-science field. Of the accepted matches, 79\% of women's and of men's pass a data check (surname in the profile name and at least two words of its main topics on the Wikipedia page).

\subsection*{Network construction}
Nodes are pages; a directed edge $A\to B$ means the wikitext of $A$ links to $B$ (redirects resolved, self-links removed). The main network has 9,680 nodes and 10,563 edges. A second, prose-only edge list (7,589 edges) excludes reference-type sections, templates including infoboxes and navigation boxes, tables and list lines. The field $\times$ gender model uses the three-field network (12,463 nodes, 12,203 edges). Networks were handled with NetworkX \cite{hagberg2008networkx}.

\subsection*{Text processing}
Plain text excludes reference-type sections, templates and tables. Person names found by spaCy's small English model were replaced by a placeholder; pronouns, gender words and titles were masked; kinship words were made neutral; names of women-only prizes and societies were masked. Tokens are lower-case alphabetic strings of three or more letters.

\subsection*{Measurements}
\emph{In-degree} counts links received from other scientists' pages; it captures how often a person appears in other biographies and misses links from non-biography pages. \emph{Out-degree} and \emph{isolates} capture whether the page itself connects; they miss why. The \emph{classifier AUC} captures whether wording differs systematically; it misses what the difference is about. \emph{Word log-odds} use the share of pages containing a word, with a Beta prior worth ten pages centred on the pooled share, an adaptation of \cite{monroe2008fightin} to document frequencies, computed within field $\times$ decade cells and combined with inverse-variance weights.

\subsection*{Validation and null models}
For RQ1, gender was permuted within strata on the fixed network (1,000 runs); because labels and not links are permuted, the test does not assume that pages' in-degrees are independent. Null models and community detection follow the course text \cite{coscia2025atlas}. Count models were fitted with statsmodels \cite{seabold2010statsmodels} (negative binomial and Poisson with robust standard errors, also clustered by Louvain community; logit plus zero-truncated negative binomial for the hurdle), with log citations, field, birth decade and deceased as controls. For the text, logistic regression (scikit-learn \cite{pedregosa2011scikit}, $C=1$, balanced classes) was scored by five-fold cross-validation and the full procedure repeated 200 times with gender permuted within field $\times$ decade. For same-gender linking, gender was permuted within Louvain community $\times$ decade for ten partitions. A check of 100 matches found 100 correct (per gender 50 of 50, 95\% interval 93--100\%); this check and the 300 labelled sentences were AI-assisted, in two rounds. In the second, matches were accepted only if the OpenAlex name or an alias contained the Wikipedia surname and the profile's main OpenAlex topics appeared on the page; the sentence labeller had seen the first-round labels, so $\kappa=0.99$ (personal) and $0.98$ (professional) are upper bounds. The sentences served to test a personal-life word list, which we then dropped: it found about a third of personal sentences and could not be shown to work equally well on women's and men's pages.

\subsection*{Code and data availability}
Code, data log, processed data and a walk-through notebook: \url{https://github.com/mahuizg/SocialGraphs-Interactions}, folder \texttt{pipeline\_revised}.

\bibliography{references}
\end{document}
